% ------------------------------------------------------------------
%  Capitolo 24 — Valutare per far imparare
%  Parte VI
%  Ricerca di riferimento: ricerca 01 §5.1, §9.5; 02 §9.3, §10; ricerca 08 (da fare)
%  Voci bibliografiche principali: yang2021, agarwal2014, roediger2011, hattie2007, kluger1996, wisniewski2020, black1998
%  Epigrafe: ✔ verificata sul testo (vedi ricerca 03 e registro, ricerca 00)
%  Scheda del capitolo: ricerca/06_indice-ragionato.md, capitolo 24
% ------------------------------------------------------------------
\chapter{Valutare per far imparare}
\label{cap:valutare}

\epigrafe[io non gli insegno nulla, ma gli domando tutto]{ἐγὼ τοῦτον οὐδὲν διδάσκω, ἀλλ᾽ ἐρωτῶ πάντα}{Platone, Menone 82e}

\iniziale{Q}{uesto capitolo} \segnaposto{apertura: da scrivere — vedi la scheda nell'indice ragionato}.

\section{L'esame come strumento di apprendimento}

\segnaposto{da scrivere}

\section{Quiz frequenti e prove in itinere}

\segnaposto{da scrivere}

\section{Esami cumulativi e sessioni}

\segnaposto{da scrivere}

\section{Ripensare l'esame orale}

\segnaposto{da scrivere}

\section{Un feedback che aiuta}

\segnaposto{da scrivere}

\begin{daricordare}
\segnaposto{il principio del capitolo in due righe}
\end{daricordare}

\begin{domande}
  \item \segnaposto{domanda di recupero su questo capitolo}
  \item \segnaposto{domanda di applicazione a una materia che studi}
  \item \segnaposto{domanda di ripresa su un capitolo precedente}
\end{domande}

\begin{sintesi}
  \item \segnaposto{punto di sintesi}
\end{sintesi}

\finecapitolo
